\section{EM 用下界交替完成隐变量学习}

隐变量模型的共同困难是：若隐变量已知，参数容易估计；若参数已知，隐变量后验容易计算；两者同时未知时，边缘 likelihood 却难直接优化。EM 把这个循环依赖变成两个可执行步骤。

先对任意分布 \(q(z)\) 写出

\[
\log p(x\mid\theta)
=\mathcal L(q,\theta)
+\KL
\bigl(q(z)\|p(z\mid x,\theta)\bigr),
\]

其中

\[
\mathcal L(q,\theta)
=\E_q[\log p(x,z\mid\theta)]
+H(q),
\]

\(H(q)=-\E_q[\log q(z)]\) 是 \(q\) 的熵。这个恒等式可以直接验证：把两项按定义展开，

\[
\mathcal L(q,\theta)
+\KL\bigl(q(z)\|p(z\mid x,\theta)\bigr)
=\E_q\left[
\log\frac{p(x,z\mid\theta)}{q(z)}
+\log\frac{q(z)}{p(z\mid x,\theta)}
\right]
=\log p(x\mid\theta),
\]

因为两个 \(\log q(z)\) 项抵消，联合密度与条件密度之比恰好留下 \(p(x\mid\theta)\)，而它对 \(z\) 是常数，期望不起作用。KL 散度非负，所以 \(\mathcal L\) 是 log-likelihood 的下界，也称 evidence lower bound（ELBO）。这个恒等式给 EM 两步以明确含义。

\subsection{E 步把下界贴紧，M 步抬高下界}

在第 \(t\) 轮，E 步令

\[
q^{(t)}(z)
=p(z\mid x,\theta^{(t)}).
\]

此时 KL 为零，下界在当前参数处恰好等于真实 log-likelihood。随后 M 步选择

\[
\theta^{(t+1)}
\in\argmax_\theta
\E_{q^{(t)}}
[\log p(x,z\mid\theta)].
\]

因为 E 步之后熵 \(H(q^{(t)})\) 对 \(\theta\) 是常数，这等价于提高同一个 ELBO。新参数处的真实 log-likelihood 又至少等于该下界，于是

\[
\ell(\theta^{(t+1)})
\ge\mathcal L(q^{(t)},\theta^{(t+1)})
\ge\mathcal L(q^{(t)},\theta^{(t)})
=\ell(\theta^{(t)}),
\]

其中 \(\ell(\theta)=\log p(x\mid\theta)\)（多样本时为各样本对数似然之和）。三步依次是：KL 非负给出下界，M 步提高下界，E 步已让下界在旧参数处贴紧。

这就是 EM 的单调性保证。它只保证似然不下降，不保证达到全局最大值，也不排除鞍点、退化解或极慢收敛。

\subsection{GMM 让两步变得具体}

对 Gaussian mixture，E 步计算上一节的 responsibility \(\gamma_{ik}\)。令

\[
N_k=\sum_i\gamma_{ik},
\]

M 步得到加权统计量：

\[
\pi_k^{\text{new}}=\frac{N_k}{n},
\]

\[
\mu_k^{\text{new}}
=\frac1{N_k}\sum_i\gamma_{ik}x_i,
\]

\[
\Sigma_k^{\text{new}}
=\frac1{N_k}
\sum_i\gamma_{ik}
(x_i-\mu_k^{\text{new}})
(x_i-\mu_k^{\text{new}})^\top.
\]

若 \(\gamma_{ik}\) 只有 0 和 1，这些式子就是按已知类别分组的普通均值与 covariance。EM 用 posterior 权重保留归属不确定性，硬猜标签再重训会丢掉这部分信息。

例如用振动与温度数据区分机台运行工况时，\(\gamma_{ik}=0.8\) 可以读作：在当前混合模型与参数下，这条记录对第 \(k\) 个成分有较高 posterior responsibility。它不能直接改写成``机台有八成概率处于某个真实故障状态''。真实工况可能连续漂移，一个偏斜工况也可能被几个 Gaussian 成分共同近似；反过来，两个物理状态若在现有传感器上重叠，也可能被合成一个成分。要赋予成分物理名称，还需要维修记录、受控切换或额外传感器等模型外证据。

\subsection{单调曲线仍可能通向坏答案}

若初始均值都落在同一密集区域，EM 可能停在较差局部最优；若某个成分抓住单个样本，似然又可能朝上一节讨论过的 covariance 塌缩退化方向上升。单调性不能告诉我们模型选择对不对，也不能告诉我们最高似然解是否有科学意义。

可靠实现通常需要：

\begin{itemize}
\tightlist
\item
  在 log 域或用 log-sum-exp 归一化 E 步；
\item
  对 covariance、权重等参数施加有效约束；
\item
  每轮记录真实 observed-data log-likelihood，而不只记录 complete-data 期望；
\item
  用多个有差异的初值拟合；
\item
  以相对似然增量、参数变化和最大轮数共同停止；
\item
  在独立数据上比较解，而不是按训练似然单独排名。
\end{itemize}

若精确 E 步不可计算，就不能再直接声称使用标准 EM。变分 EM 用受限的 \(q\) 近似后验，Monte Carlo EM 用样本近似期望；generalized EM 则不完全最大化 M 步，只要求足够提高目标。每种变体都需要重新说明单调性或近似误差的条件。
